\section{The trace sensor and its weighted Ward limit}
\label{sec:trace-prosper-sensor}

Fix one selected family and a gauge group given as a closed subgroup of a
finite-dimensional unitary group, whose real Lie algebra is nonabelian and
simple. Let $E$ be the real
finite-dimensional space of grade-five vector coefficient arrays obeying the
reality and signed covariance conditions, and let $H$ be the corresponding
reference vector table space. The free projection $f:E\to H$ applies the
reference field projection in each vector component. Let $x\in E$ be the
coefficient array of the free equation of motion. The free rank calculation
in Section~\ref{sec:ex-lattice-trace-free-rank} gives $x\ne0$ and
$\ker f=\mathbb R x$. To fix the free Ward normalization, write the raw
equation-of-motion array as $J^{(3)}_b=\sum_{i,a}\langle F_{bi},\partial_a
F_{ai}\rangle$, and let
\[
 (J^{(1)})_b=\sum_i\langle F_{bi},\partial_b F_{bi}\rangle,\qquad
 (J^{(2)})_b=\sum_{i,j}\langle F_{ij},\partial_b F_{ij}\rangle.
\]
Here $F$ and $\partial F$ are the free curvature jet and its first derivative,
and the inner product is the gauge-invariant real color inner product. Each
array is real and invariant under the color action. For a coordinate
reflection, every contracted index contributes its sign twice, leaving only
the sign of the output index $b$; permutations merely relabel the sums. Thus
each array has the required signed-vector covariance. These three arrays form
a basis of the raw vector space. To see independence, choose
a unit color vector and evaluate the zeroth component on three jets. In each,
the only nonzero curvature pair, up to antisymmetry, is respectively $F_{12}$,
$F_{01}$, and $F_{01}$. The nonzero derivative entries, again up to
antisymmetry, are respectively
\[
 \begin{array}{c|l}
 1&\partial_0F_{12}=1,\quad \partial_1F_{20}=-\tfrac12,\quad
       \partial_2F_{01}=-\tfrac12,\\
 2&\partial_0F_{01}=1,\quad \partial_2F_{21}=-1,\\
 3&\partial_0F_{01}=1.
 \end{array}
\]
Each scalar stands for that scalar times the chosen color vector. These jets
satisfy the Bianchi identity; for example a curvature jet $(F_{ab},K_{cab})$
with $K_{cab}+K_{abc}+K_{bca}=0$ is realized by the polynomial potential
$A_b(z)=\tfrac12z^aF_{ab}+\tfrac13z^az^cK_{cab}$. It is in radial gauge,
since $z^bA_b(z)=0$ by antisymmetry in $a,b$. Its curl at the origin
is $F_{ab}$, and its first curl derivative is
\[
 \partial_c(\partial_aA_b-\partial_bA_a)(0)
 =\tfrac13(2K_{cab}-K_{abc}-K_{bca})=K_{cab},
\]
where the last equality uses the cyclic identity. Direct substitution gives
the evaluation matrix
\[
 \begin{pmatrix}0&2&0\\1&2&0\\1&2&1\end{pmatrix}
 \quad\text{for }(J^{(1)}_0,J^{(2)}_0,J^{(3)}_0).
\]
Its determinant is nonzero. The free-rank calculation gives the matching
upper bound three for the raw vector dimension, so the arrays form a basis.
The label-coordinate identification carries $x$ to $J^{(3)}$. Every raw vector
therefore has unique real coordinates $(w_0,w_1,w_2)$ in the form
\[
 2(w_0-w_2)J^{(1)}+(2w_1+\tfrac12w_2)J^{(2)}+w_2J^{(3)}.
\]
Indeed, the $J^{(3)}$ coefficient determines $w_2$, and the other two
coefficients then determine $w_0$ and $w_1$.
To see the Ward normalization, let $\mathcal I=\sum_{i,j}\langle F_{ij},F_{ij}\rangle$
and consider the centered stress tensor with diagonal entries
$T_{aa}=w_0\sum_i\langle F_{ai},F_{ai}\rangle+w_1\mathcal I$ and off-diagonal
entries $T_{ab}=w_2\sum_i\langle F_{ai},F_{bi}\rangle$. Its divergence
has the displayed raw-vector coordinates. Indeed the diagonal derivatives
contribute $2w_0J^{(1)}+2w_1J^{(2)}$. For the off-diagonal term, first
sum over all $a$. One product-rule term is $J^{(3)}$. For the other, write
$B=\sum_{a,i}\langle F_{ai},\partial_aF_{bi}\rangle$. The cyclic Bianchi
identity, antisymmetry, and interchange of $a,i$ give
\[
 B=\sum_{a,i}\langle F_{ai},\partial_bF_{ai}\rangle
    -\sum_{a,i}\langle F_{ai},\partial_iF_{ab}\rangle
  =J^{(2)}-B.
\]
Hence the two terms sum to $J^{(3)}+\tfrac12J^{(2)}$. Removing $a=b$
subtracts $2J^{(1)}$.
Its trace is $(w_0+4w_1)\mathcal I$. The scalar normalization is $\mathcal I/4$, and the
free Ward functional $\rho:E\to\mathbb R$ is minus the coefficient of
that scalar label in the trace. The label is nonzero, as seen by taking
$F_{01}$ to be a unit color vector and the other independent curvature
components to be zero. Hence $\rho=-4(w_0+4w_1)$ on the above coordinates. In
particular $J^{(3)}$ has coordinates $(1,-\tfrac14,1)$ and $\rho(x)=0$.

Let $k_i$ be the selected initial levels, with $k_i\to\infty$. Fix the
reference level $j_0$ of the grade-five coefficient space. All finite-volume
maps and estimates below are used at levels $j\ge j_0$; in particular,
$k_i\ge j_0$ eventually. The running rate is
$\widehat\beta_n=\bar\beta+\gamma\ell_n$, where $\gamma>0$,
$\widehat\beta_n>1$, and the schedule scale $\ell_n\to\infty$. For a level
$j$, let $q(j)>j$ be the first level whose logarithmic scale separation from
$j$ is at least $(\log\widehat\beta_j)^{48}$, and put $n_i=q(k_i)$.
Consequently $n_i\to\infty$,
$\widehat\beta_{n_i}\to\infty$, and the positive weight
\begin{equation}
 \mathfrak a_i=\exp\!\left\{\frac15\bigl(\log\log\widehat\beta_{n_i}\bigr)^2\right\}
 \longrightarrow\infty.
 \label{eq:trace-sensor-weight}
\end{equation}
This is the $1/5$ age weight at the \emph{postselected} level $n_i$;
it is not the reciprocal chart scale.

We shall use the following clock and arch inputs. The schedule has natural
block sizes $L_j$ and real increments $\ell_{j+1}-\ell_j=\log L_j$, with
$L_j\ge3$ and, eventually,
$L_j\le\log\widehat\beta_j$. The clock offset $\bar\beta$ and the arch
scales $b_j$ are positive, and for every $\varepsilon>0$,
\begin{equation}
 |b_j-\widehat\beta_j|\le\varepsilon\ell_j
 \quad\text{for all sufficiently large }j.
 \label{eq:trace-sensor-clock-input}
\end{equation}
These are selected-run inputs.

\begin{lemma}[Arch-scale ratio]\label{lem:trace-sensor-arch-ratio}
Under the clock inputs above, $b_{k_i}/b_{n_i}\to1$.
\end{lemma}
\nopagebreak
\begin{proof}
Put
$\chi=\widehat\beta_k$ and $g=(\log \chi)^{48}$, and let $q(k)=m+1$.
By minimality, $\ell_m-\ell_k\le g$: if $m=k$ it is zero, and otherwise
the preceding level has not reached the threshold. For large $k$,
$\gamma g\le \chi$, so $\widehat\beta_m\le2\chi$. The last increment obeys
\[
 \ell_{m+1}-\ell_m=\log L_m\le L_m
 \le\log\widehat\beta_m\le\log(2\chi)\le g.
\]
Consequently $0\le\ell_{q(k)}-\ell_k\le2g$, and therefore
\[
 \frac{\widehat\beta_{q(k)}}{\widehat\beta_k}
 =1+\gamma\frac{\ell_{q(k)}-\ell_k}{\widehat\beta_k}
 \longrightarrow1,
\]
because $(\log \chi)^{48}/\chi\to0$. Also
$\widehat\beta_j=\bar\beta+\gamma\ell_j\ge\gamma\ell_j$, so
\eqref{eq:trace-sensor-clock-input} gives $b_j/\widehat\beta_j\to1$.
Since both $k_i$ and $n_i=q(k_i)$ tend to infinity,
\begin{equation}
 \frac{b_{k_i}}{b_{n_i}}\longrightarrow1.
 \label{eq:trace-sensor-arch-ratio}
\end{equation}
\end{proof}

Write $s_i:E\to H$ for the selected, normalized grade-five chart map. For
one vector component let $\mathsf V$ be the complex coefficient space,
$\mathcal F$ the reference field space and $S_5$ the selected
interacting field space. The interacting projection is $P:\mathsf V\to
S_5$, the scale-$i$ identification is a linear isomorphism
$X_i:\mathcal F\to S_5$, and the selected terminal lift is
$Y_i:\mathcal F\to\mathsf V$. Write $M_i:\mathsf V\to\mathsf V$
for the invertible arch map at $n_i$, $t_i>0$ for the threshold there, and
$\operatorname{core}:\mathsf V\to\mathsf V$ for the projection deleting
the unit coordinate: if $e_0$ is the unit coordinate vector, then
$\operatorname{core}v=v-v_0e_0$. These component maps are complex-linear; $s_i$ and $u_i$ act real-linearly on the
corresponding covariant arrays. With arch scale $b_n>0$ and dilation radius $r_i>0$, set
$c_i=b_{n_i}r_i^5$. On an eventual common tail, for each vector component
indexed by $\alpha$,
\begin{equation}
 (s_i y)_\alpha=X_i^{-1}P(c_i t_i^2 M_i^{-1}y_\alpha),\qquad
 (u_i h)_\alpha=c_i^{-1}t_i^{-2}\operatorname{core} M_iY_i h_\alpha.
 \label{eq:trace-sensor-selected-formulas}
\end{equation}
The second formula is for the selected lift. Its well-definedness as a map
into $E$ follows from the chart symmetries below. On the finitely many indices
outside this tail, choose arbitrary real-linear maps $u_i:H\to E$. The limiting right inverse
$U$ applies the fixed reference lift to each component. For that
lift, the reference field is the quotient of the coefficient space by its
free equation subspace $K$. For the finite signed-permutation group $\Gamma$,
write $g z$ for the action on $z\in\mathsf V$ and equip $\mathsf V$
with the Hermitian form
\[
 \langle z,w\rangle_*=\sum_{g\in\Gamma}\sum_a
             \overline{(gz)_a}(gw)_a.
\]
It is positive definite because the identity term occurs. Reindexing the
sum proves invariance under $\Gamma$; conjugation is antiunitary because
it commutes with this action. The free equations are real and signed
permutation equivariant; their vanishing subspace $K$ is therefore stable
under these symmetries. Each quotient class has a
unique representative in $K^\perp$, and taking that representative is a
linear section intertwining the signed action and conjugation. Applied
componentwise, it therefore maps $H$ into $E$ and satisfies $fU=\mathrm{id}_H$.
These formulas use the same selected history and the same index $i$; the
symbols $n_i$ and $b_n$ are as above (with $b_n$ denoting the arch scale).
The following identities hold on a common tail. The first was just proved;
the second is a terminal projection hypothesis, and the third follows from
the chart symmetries as shown below:
\begin{equation}
 fU=\mathrm{id}_H,\qquad P Y_i=X_i,\qquad
 \operatorname{core} M_iY_i h_\alpha=M_iY_i h_\alpha.
 \label{eq:trace-sensor-chart}
\end{equation}
The last equality holds simultaneously for every $h\in H$ and every component
$\alpha$. On the same tail, $Y_i$ and $M_i$ intertwine signed permutations
and conjugation. Thus $M_iY_i$ carries real covariant arrays to real
covariant arrays. Inversion of all four axes acts by $-1$ on each lifted
component and fixes the unit coefficient, so that coefficient vanishes; the
core deletes only that coefficient. This proves the last identity in
\eqref{eq:trace-sensor-chart}. The scalar $c_i^{-1}t_i^{-2}$ is real, so the
second formula in \eqref{eq:trace-sensor-selected-formulas} also defines a
real-linear map $u_i:H\to E$. The chart identities imply, on this same tail,
\begin{align*}
 (s_i u_i h)_\alpha
 &=X_i^{-1}P\!\left(c_it_i^2M_i^{-1}
       (c_i^{-1}t_i^{-2}\operatorname{core}M_iY_i h_\alpha)\right)\\
 &=X_i^{-1}P(Y_i h_\alpha)=h_\alpha.
\end{align*}
Thus $s_i u_i=\mathrm{id}_H$ eventually, uniformly over $h$. The construction
of the chart and its symmetry and terminal projection properties are assumed
here, not consequences of the limiting rank alone.

We also record the quantitative inputs behind convergence of the selected
right inverses. For any level $j$, write $\vartheta_j>0$ for its threshold,
$M(j)$ for its arch map, and $W_j=\vartheta_j^{-2}\operatorname{core}M(j)$; thus $M_i=M(n_i)$
and $t_i=\vartheta_{n_i}$. Let $|\cdot|_1$ be the sum of absolute values of the
finite coefficient coordinates, and set
\[
 D_i=b_{k_i}r_i^5,\quad
 \eta_i=\exp\!\left\{-\tfrac12(\log\log\widehat\beta_{k_i})^2\right\},
 \quad
 \delta_i=\exp\!\left\{-\tfrac25(\log\log\widehat\beta_{k_i})^2\right\}.
\]
Both $\eta_i$ and $\delta_i$ tend to zero. For fixed $h\in H$ and
component $\alpha$, write $y_i=Y_i h_\alpha$ and $p=(Uh)_\alpha$.
The next two estimates follow from the finite-volume inputs listed in
(S1)--(S6) below; they are derived consequences, not new sensor hypotheses.

\begin{lemma}[Initial-level price]\label{lem:trace-sensor-initial-price}
Under those inputs, for each fixed $h$ and $\alpha$, eventually
\begin{equation}
 |W_{k_i}y_i-D_ip|_1\le 2\eta_iD_i|p|_1.
 \label{eq:trace-sensor-initial-price}
\end{equation}
\end{lemma}
\nopagebreak
\begin{proof}
The finite-volume indices are pairs $\omega=(N,\nu)\in\mathbb N^2$,
where $N$ is the mesh cutoff index and $\nu$ the period index. Fix the selected
ultrafilter $\mathcal L$ refining the fused order of these pairs. Write
$\mathcal F_{\mathrm{fuse}}$ for its fused-order filter; the refinement is
$\mathcal L\le\mathcal F_{\mathrm{fuse}}$ in the filter order; it
contains both coordinate tails, which are the consequences used below. The initial-level
estimate follows from a finite-volume estimate on this ultrafilter. For each
sufficiently large $i$, fixed $h$ and component $\alpha$, use the same $y_i$
and $p$ as above. Suppose finite-volume coefficients $Z_i(\omega)$ and normalizations
$B_i(\omega)$, positive on an ultrafilter-large set, satisfy
\begin{equation}
 |Z_i(\omega)-B_i(\omega)p|_1\le
 2\eta_i B_i(\omega)|p|_1
 \quad\text{on an ultrafilter-large set},
 \label{eq:trace-sensor-finite-price}
\end{equation}
and converge along that ultrafilter to $W_{k_i}y_i$ and $D_i$,
respectively. These limits have a common source. Let $\operatorname{Col}_j(\omega)$
be the finite-volume transport at level $j$. Assume that for each fixed
admissible $j$, whenever $v(\omega)\to v$ along the selected ultrafilter,
$\operatorname{Col}_j(\omega)v(\omega)\to M(j)v$; also assume the finite-volume
running parameter at level $j$ tends to $b_j$. Write $\beta_i(\omega)$
for that parameter at $k_i$. The selected finite-volume lift satisfies $\lambda_i(\omega)\to y_i$. The definitions give
\[
 Z_i(\omega)=\vartheta_{k_i}^{-2}\operatorname{core}\operatorname{Col}_{k_i}(\omega)\lambda_i(\omega),
 \qquad B_i(\omega)=\beta_i(\omega)r_i^5.
\]
The two stated limits follow by continuity of $\operatorname{core}$, with $i$ fixed.
The joint transport and lift convergence remain hypotheses; the
finite-volume price follows from the tested-map estimates below.
Continuity of the coefficient norm and the closedness of the inequality
give \eqref{eq:trace-sensor-initial-price}. The ultrafilter limit is taken
for each fixed $i$ before the limit $i\to\infty$.

The finite-volume price has a useful decomposition into its tested-map and
calibration inputs. Fix $i$ and $\omega=(N,\nu)$ with $N>k_i$; suppress both
indices, and let $\mathsf A$ and $\mathsf T$ be the transported reference
map and its tested approximation. They are complex-linear endomorphisms $\mathsf A,\mathsf T:\mathsf V\to\mathsf V$ of the finite
complex coefficient space and send the reference unit $\mathbf1$ to the
unit coordinate vector $e_0$. Assume that $\mathsf A$ is injective and, for
every input $v$,
\begin{equation}
 |\mathsf A v-\mathsf T v|_1\le\tfrac12|\mathsf A v|_1,
 \qquad
 |\mathsf A v-\mathsf T v|_1\le\eta_i|\mathsf A v|_1.
 \label{eq:trace-sensor-tested-price}
\end{equation}
These are separate tested-map and relative-price estimates. The first also
shows that $\mathsf T$ is injective: if $\mathsf T v=0$, then
$|\mathsf A v|_1\le\tfrac12|\mathsf A v|_1$, so $\mathsf A v=0$
and $v=0$. Since $\mathsf V$ is finite-dimensional, $\mathsf T$ is invertible. For a given $v$, put
$c=(\mathsf T v)_0\in\mathbb C$ and $v^\circ=v-c\mathbf1$. Then
$\mathsf T v^\circ=\operatorname{core}\mathsf T v$ and
$\operatorname{core}\mathsf A v^\circ=\operatorname{core}\mathsf A v$. The first bound in
\eqref{eq:trace-sensor-tested-price} and the triangle inequality give
$|\mathsf A v^\circ|_1\le2|\mathsf T v^\circ|_1$. Applying $\operatorname{core}$ to
the second bound, and using its nonexpansiveness, yields
\begin{equation}
 |\operatorname{core}\mathsf A v-\operatorname{core}\mathsf T v|_1
 \le2\eta_i|\operatorname{core}\mathsf T v|_1.
 \label{eq:trace-sensor-centered-price}
\end{equation}

For the particular selected lift $v=\lambda_i(\omega)$, its definition
identifies the tested image and finite-volume normalization. Write
$d_{k_i}>0$ for the lattice spacing at $k_i$, put
$\sigma_i=r_i/d_{k_i}$, and let $\beta_i(\omega)>0$ be the finite-volume
running parameter. The transported reference map is
$\mathsf A=d_{k_i}^{-5}\operatorname{Col}_{k_i}(\omega)$, which explains the
normalization of $Z_i$ below. The vector covariance of $h$ makes each $h_\alpha$ odd
under total inversion; the reference lift intertwines that action, so
$p=(Uh)_\alpha$ is odd as well. We recall the short degree argument for
the reference frame. Its basis elements are homogeneous in position degree
and factor degree, and the only factor-degree-zero basis element is the unit. Each
jet factor has position weight at least two; since the frame has degree at
most five, a nonunit basis element has at most two factors. A one-factor
invariant vanishes by the invariant-linear-form argument in the low-grade
label calculation of Section~\ref{sec:ex-lattice-trace-free-rank}. Thus every nonunit basis element has exactly two
factors and position degree four or five. Total inversion acts on a homogeneous
position-degree-$d$ label by $(-1)^d$: a jet of order $r$ has both inversion
sign $(-1)^{r+2}$ and position weight $r+2$. Consequently every nonzero
nonunit coordinate of the odd vector $p$ has position degree five and factor
degree two. The identities used here are
\begin{equation}
 \operatorname{core}\mathsf T\lambda_i=\beta_i\vartheta_{k_i}^2\sigma_i^5p,\qquad
 Z_i=\vartheta_{k_i}^{-2}d_{k_i}^5\operatorname{core}\mathsf A\lambda_i,\qquad
 B_i=\beta_i r_i^5.
 \label{eq:trace-sensor-finite-calibration}
\end{equation}
The first identity follows from the definition of the selected finite-volume
lift and the grade-five scaling rule. Put
$\xi=(\sqrt{\beta_i})^{-1}/\vartheta_{k_i}$. For a frame coordinate
$\iota$ with position degree $d_\iota$ and factor degree $n_\iota$,
define the position and factor dilations by
\[
 (\operatorname{Pos}_t v)_\iota=t^{d_\iota}v_\iota,\qquad
 (\operatorname{Dil}_t v)_\iota=t^{n_\iota}v_\iota.
\]
For $t>0$, these coordinate formulas give
$\operatorname{Dil}_t^{-1}=\operatorname{Dil}_{t^{-1}}$. On the selected
tail, where $i\ge5$, the finite-volume lift is
\[
 \lambda_i=\mathsf T^{-1}\operatorname{Dil}_{\xi^{-1}}
                    \operatorname{Pos}_{\sigma_i}p.
\]
Consequently $\operatorname{Dil}_\xi\mathsf T\lambda_i=\operatorname{Pos}_{\sigma_i}p$.
The reference
section $p$ is odd and has nonunit degree five, so this weight is
$\sigma_i^5p$. On these nonunit coordinates the inverse coefficient
dilation $\operatorname{Dil}_{\xi^{-1}}$ acts by $\xi^{-2}$. Applying that inverse
and $\operatorname{core}$, with $\operatorname{core}p=p$, gives the first identity in
\eqref{eq:trace-sensor-finite-calibration}, because
$\xi^{-2}=\beta_i\vartheta_{k_i}^2$. This is an exact conversion of the
lift and grading normalizations.
The other two identities specify the actual normalization. Multiplying
\eqref{eq:trace-sensor-centered-price} by
$\vartheta_{k_i}^{-2}d_{k_i}^5$ and using $r_i=d_{k_i}\sigma_i$ gives exactly
\eqref{eq:trace-sensor-finite-price}. Thus the tested-map estimates,
selected-lift construction and finite-volume limits are the remaining
inputs for the initial-level price, rather than an additional assumed
sensor bound.
\end{proof}

\begin{lemma}[Transport price]\label{lem:trace-sensor-transport-price}
Under (S1)--(S6), for each fixed $h$ and $\alpha$, eventually
\begin{equation}
 |W_{k_i}y_i-W_{n_i}y_i|_1\le\delta_i|W_{n_i}y_i|_1.
 \label{eq:trace-sensor-transport-price}
\end{equation}
\end{lemma}
\nopagebreak
\begin{proof}
For the one-step input, fix $i$, $h$ and $\alpha$ and write
$N_i=n_i-k_i>0$. On a set belonging to the selected ultrafilter, suppose
there are finite-volume coefficient-vector functions $z_j(\omega)$ for
$k_i\le j\le n_i$, all formed from the same terminal label $y_i$, with
\begin{equation}
 |z_j(\omega)-z_{j+1}(\omega)|_1\le
 \frac{\delta_i}{2N_i}|z_{j+1}(\omega)|_1
 \quad(k_i\le j<n_i).
 \label{eq:trace-sensor-step-input}
\end{equation}
Here $z_j(\omega)=\vartheta_j^{-2}\operatorname{core}\operatorname{Col}_j(\omega)y_i$. For each
fixed $j$ in this finite interval, the joint transport input above, applied
to the constant label $y_i$, gives the ultrafilter limit $W_jy_i$. We derive
the one-step price below from the underlying residual and budget inputs, and
now prove
\eqref{eq:trace-sensor-transport-price} from it.

Fix $\omega$ in that set and abbreviate $z_j=z_j(\omega)$. Put
$a=\delta_i/N_i$. For large $i$ we have $\delta_i\le1$. Backward
induction gives, for every $j$ in the interval,
\[
 |z_j-z_{n_i}|_1\le a(n_i-j)|z_{n_i}|_1.
\]
It is immediate at $j=n_i$. At the induction step, the bound at $j+1$
and $aN_i\le1$ give $|z_{j+1}|_1\le2|z_{n_i}|_1$. The one-step input
therefore bounds $|z_j-z_{j+1}|_1$ by $a|z_{n_i}|_1$, and the triangle
inequality proves the next step. At $j=k_i$ this is
$|z_{k_i}-z_{n_i}|_1\le\delta_i|z_{n_i}|_1$. Passing to the ultrafilter
limit in the finite-dimensional coefficient space yields
\eqref{eq:trace-sensor-transport-price}. The finite intersection over
the interval is taken before this limit, for each fixed $i$; no uniform
limit in $i$ or new choice of history is used.

The one-step price follows from the following assumptions. Write
$\widetilde z_j=\operatorname{Col}_j(\omega)y_i$ for the finite-volume
coefficient before threshold normalization; thus
$z_j=\vartheta_j^{-2}\operatorname{core}\widetilde z_j$. Let the
deterministic block transport be
$(\operatorname{Block}_j v)_\iota=L_j^{5-d_\iota}v_\iota$.
For $a\in\mathbb R$ and level $j$, put
\[
 A_a(j)=\exp\!\left\{a(\log\log\widehat\beta_j)^2\right\},\qquad
 g_j=\frac{2(1+C_{j+1})}{A_{1/2}(j+1)},
\]
where $C_j\ge0$ is the selected real column budget. The budget input is
$C_j\le\log\widehat\beta_j$ eventually. For all levels in the interval,
on the same ultrafilter-large set, assume
\begin{align}
 \frac{\vartheta_{j+1}}{\vartheta_j}&\le2, &
 |\widetilde z_j-\operatorname{Dil}_{\vartheta_j/\vartheta_{j+1}}\operatorname{Block}_j\widetilde z_{j+1}|_1
 &\le g_j|\operatorname{core}\widetilde z_{j+1}|_1.
 \label{eq:trace-sensor-raw-step}
\end{align}
In the selected construction, for each fixed level $j$ the residual and
threshold bounds hold for all sufficiently large $\nu$, uniformly in $N>j$
and the terminal label. For fixed $i$, take the maximum
of the finitely many period thresholds for $k_i\le j<n_i$ and also require
$N\ge n_i$; their intersection belongs to the same ultrafilter. This
justifies using one large set for the entire finite interval, without a
uniform choice over $i$.

The terminal label $y_i$ is
odd under inversion of all axes, and the finite-volume transport preserves
that inversion; hence every $\widetilde z_j$ is odd. The reference-frame classification
used above therefore shows that a nonzero nonunit coordinate of $\widetilde z_{j+1}$ has
position degree five and factor degree two. On a coordinate of position
degree $d$, block transport multiplies by
$L_j^{5-d}$; on one of factor degree $n$, coefficient dilation by $t$
multiplies by $t^n$. Thus, on every nonunit coordinate of $\widetilde z_{j+1}$,
the block factor is one and the dilation factor is
$(\vartheta_j/\vartheta_{j+1})^2$. The core deletes the unit coordinate, so
\begin{equation}
 \operatorname{core}\operatorname{Dil}_{\vartheta_j/\vartheta_{j+1}}\operatorname{Block}_j\widetilde z_{j+1}
 =\left(\frac{\vartheta_j}{\vartheta_{j+1}}\right)^2\operatorname{core}\widetilde z_{j+1}.
 \label{eq:trace-sensor-grade-scaling}
\end{equation}
The core projection is nonexpansive in $|\cdot|_1$.
Applying $\operatorname{core}$ to the residual and dividing by $\vartheta_j^2$ gives
\[
 |z_j-z_{j+1}|_1
 \le g_j\left(\frac{\vartheta_{j+1}}{\vartheta_j}\right)^2|z_{j+1}|_1
 \le4g_j|z_{j+1}|_1.
\]
We verify that
$4g_j\le\delta_i/(2N_i)$ uniformly for $k_i\le j<n_i$ when $i$ is large.
Put $\chi=\widehat\beta_{k_i}$. The clock estimates already proved imply
$\widehat\beta_{n_i}\le2\chi$ eventually. Since each increment
$\log L_j\ge1$, the overshoot bound gives $N_i\le2(\log \chi)^{48}$.
For large $i$ we have $\chi\ge e$, so
$t\mapsto\exp\{\tfrac12(\log\log t)^2\}$ is increasing for $t\ge\chi$.
Monotonicity of the clock then gives $C_{j+1}\le2\log \chi$ and
$A_{1/2}(j+1)\ge A_{1/2}(k_i)$ for large $i$. Therefore
\[
 2N_i(4g_j)\le
 \frac{32(\log \chi)^{48}(1+2\log \chi)}{A_{1/2}(k_i)}
 \le A_{-2/5}(k_i)=\delta_i.
\]
The last inequality follows because every fixed power of $\log \chi$ is
eventually bounded by $A_{1/10}(k_i)$: after taking logarithms this is
linear versus quadratic growth in $\log\log \chi$. This proves
\eqref{eq:trace-sensor-step-input}. The residual, preservation of oddness, threshold and column-budget bounds
remain hypotheses; the degree classification, grade-scaling identity,
price calculation and telescoping above prove their consequence.
\end{proof}

\begin{lemma}[Convergence of the selected lifts]\label{lem:trace-sensor-lift-limit}
Under (S1)--(S6), $u_i\to U$ in operator norm.
\end{lemma}
\nopagebreak
\begin{proof}
The preceding estimates prove $u_i h\to Uh$. Indeed put
$\mathsf L_i=D_i^{-1}W_{k_i}y_i$ and $\mathsf R_i=D_i^{-1}W_{n_i}y_i$.
The first estimate gives $|\mathsf L_i-p|_1\le2\eta_i|p|_1$, and the second
gives $|\mathsf L_i-\mathsf R_i|_1\le\delta_i|\mathsf R_i|_1$. For $\delta_i<1$, the triangle
inequality gives $|\mathsf R_i|_1\le(1+2\eta_i)(1-\delta_i)^{-1}|p|_1$ and hence
\[
 |\mathsf R_i-p|_1\le\frac{2\eta_i+\delta_i}{1-\delta_i}|p|_1.
\]
Since $c_i=b_{n_i}r_i^5$, the formula for $u_i$ gives
$(u_i h)_\alpha=(b_{k_i}/b_{n_i})\mathsf R_i$. Thus the component error is at
most
\[
 \left(\left|\frac{b_{k_i}}{b_{n_i}}-1\right|
 +\frac{b_{k_i}}{b_{n_i}}
   \frac{2\eta_i+\delta_i}{1-\delta_i}\right)|p|_1,
\]
which tends to zero. There are finitely many components, so $u_i h\to Uh$
in $E$; finite dimensionality upgrades this pointwise convergence to
operator-norm convergence.
\end{proof}

We collect the inputs used in the preceding lemmas and the proposition.
They concern one selected family, one chart and one ultrafilter throughout;
none is a new assumption of a sensor estimate.
\begin{enumerate}
\item[(S1)] The selected levels $k_i\to\infty$ and the schedule has
  $\widehat\beta_j=\bar\beta+\gamma\ell_j>1$, with $\bar\beta,\gamma>0$
  and $\ell_j\to\infty$. Its increments satisfy
  $\ell_{j+1}-\ell_j=\log L_j$, with $L_j\ge3$ and eventually
  $L_j\le\log\widehat\beta_j$. The arch scales $b_j$ are positive and,
  for every $\varepsilon>0$, satisfy
  $|b_j-\widehat\beta_j|\le\varepsilon\ell_j$ eventually in $j$.
  The postselected level $n_i=q(k_i)>k_i$ is the first level at logarithmic
  separation at least $(\log\widehat\beta_{k_i})^{48}$, and the weight is
  \eqref{eq:trace-sensor-weight}.
\item[(S2)] The finite-dimensional real spaces $E,H$, projection $f$,
  nonzero $x$, $\ker f=\mathbb Rx$ and Ward form $\rho$ are the free data
  fixed above. The chart maps $s_i$, $X_i$, $Y_i$, $M_i$ and positive scales
  use the same selected family and satisfy
  \eqref{eq:trace-sensor-selected-formulas} and $PY_i=X_i$ on a common
  eventual tail for every field. The chart maps intertwine the signed
  permutations and conjugation there; the reference section $U$ has
  $fU=\mathrm{id}_H$.
\item[(S3)] The selected ultrafilter $\mathcal L$ refines the fused order.
  For each fixed admissible $j$, $\operatorname{Col}_j(\omega)v(\omega)
  \to M(j)v$ whenever $v(\omega)\to v$ along $\mathcal L$; the finite
  running parameter tends to $b_j$. For each fixed sufficiently large $i$
  and each field component, the selected finite-volume lift
  $\lambda_i(\omega)$ converges to $Y_i h_\alpha$ along this same
  ultrafilter. All these volume limits are taken before $i\to\infty$.
\item[(S4)] On the selected tail, for every period and every $N>k_i$,
  the complex-linear endomorphisms $\mathsf A,\mathsf T$ of $\mathsf V$
  send the reference unit to $e_0$, $\mathsf A$ is injective, and the two
  bounds \eqref{eq:trace-sensor-tested-price} hold for every coefficient
  vector. For each fixed $h$ and component, the selected lift is
  \[\lambda_i=\mathsf T^{-1}\operatorname{Dil}_{\xi^{-1}}
  \operatorname{Pos}_{\sigma_i}(Uh)_\alpha,\qquad
  \xi=(\sqrt{\beta_i})^{-1}/\vartheta_{k_i}.\]
  Further,
  $\mathsf A=d_{k_i}^{-5}\operatorname{Col}_{k_i}$, with positive
  $\beta_i$, $d_{k_i}$, $r_i$, and $\sigma_i=r_i/d_{k_i}$. The
  calibration identities \eqref{eq:trace-sensor-finite-calibration} are
  derived from these data and the grading, as shown in the proof.
\item[(S5)] For each fixed level $j\ge j_0$, eventually in the period
  index $\nu$, uniformly for every $N>j$ and every terminal coefficient
  label, the threshold ratio and residual bounds
  \eqref{eq:trace-sensor-raw-step} hold with
  $\widetilde z_j=\operatorname{Col}_j(\omega)v$ for that label $v$.
  For each fixed $i$ their finite
  intersection over $k_i\le j<n_i$, together with $N\ge n_i$, is taken
  along the same fused-order ultrafilter; no tail uniform in $i$ is assumed.
\item[(S6)] The selected column budgets obey $C_j\le\log\widehat\beta_j$
  eventually. The reference lift is odd under total inversion and the
  finite-volume transport preserves oddness; the grade-five frame and
  block/dilation rules are those used above. These inputs purchase the
  transport estimate, rather than assuming it.
\item[(S7)] The marked sequence $v_i\in E$ has the three limits in
  \eqref{eq:trace-sensor-marked} below. It is arbitrary subject to those
  limits; it need not be the later medium example.
\end{enumerate}

\begin{proposition}[Selected trace sensor]
\label{prop:trace-prosper-sensor}
Assume (S1)--(S6). By Lemmas~\ref{lem:trace-sensor-arch-ratio},
\ref{lem:trace-sensor-initial-price}, \ref{lem:trace-sensor-transport-price}
and \ref{lem:trace-sensor-lift-limit}, the stated chart has an eventual
right inverse and $u_i\to U$ in operator norm.
Suppose a sequence $v_i\in E$
satisfies, in the indicated norms,
\begin{equation}
 v_i\longrightarrow x,\qquad
 \mathfrak a_i s_i(v_i)\longrightarrow0\ \text{in }H,\qquad
 \mathfrak a_i\rho(v_i)\longrightarrow0\ \text{in }\mathbb R.
 \label{eq:trace-sensor-marked}
\end{equation}
Then $s_i\to f$ in operator norm. Moreover there are linear functionals
$\mu_i:H\to\mathbb R$ and $C\ge0$ such that, for all sufficiently large
$i$, $|\mu_i(h)|\le C\|h\|$ for every $h\in H$, and for every
$\varepsilon>0$, for all sufficiently large $i$,
\begin{equation}
 |\mathfrak a_i(\mu_i(s_i y)-\rho(y))|\le\varepsilon\|y\|
 \quad\text{for every }y\in E.
 \label{eq:trace-sensor-conclusion}
\end{equation}
The tail in the second assertion may depend on $\varepsilon$, but is uniform
in $y$.
\end{proposition}
\begin{proof}
Since $\mathfrak a_i>0$ and $\mathfrak a_i\to\infty$, the first weighted limit in
\eqref{eq:trace-sensor-marked} implies $s_i(v_i)\to0$. Choose a linear
functional $A:E\to\mathbb R$ for which
$y=Uf(y)+A(y)x$ for every $y$: indeed $y-Uf(y)$ lies in $\ker f$, and its
coefficient along the nonzero $x$ is linear. Define
$T_i=u_i f+A\otimes v_i$. The chart and marked-vector limits imply
$T_i\to\mathrm{id}_E$ in operator norm. For sufficiently large $i$, $T_i$
is invertible and $T_i^{-1}\to\mathrm{id}_E$, by the geometric series for
$\mathrm{id}_E-T_i$.

On the common right-inverse tail,
\[
 s_iT_i y=f(y)+A(y)s_i(v_i).
\]
This also gives a bound without assuming the norms of $s_i$ in advance.
Put $\epsilon_i=\|T_i-\mathrm{id}_E\|$; for $\epsilon_i<1$,
\[
 \|s_i-f\|\le
 \frac{\|A\|\,\|s_i(v_i)\|+\|f\|\epsilon_i}{1-\epsilon_i}.
\]
Indeed subtract $fT_i$ from the displayed identity and multiply on the
right by $T_i^{-1}$. Both terms in the numerator tend to zero, proving
$s_i\to f$. Now apply the weighted rank-one
receiver lemma (Lemma~\ref{lem:ex-lattice-trace-weighted-price}) with
$(v_i,u_i,w_i,e_i)$ there equal to $(s_i,u_i,v_i,\mathfrak a_i)$ here. Its two
weighted hypotheses are precisely the last two limits in
\eqref{eq:trace-sensor-marked}; its other hypotheses follow from the
preceding paragraph and \eqref{eq:trace-sensor-chart}. It gives
$\mu_i=\rho u_i$ with bounded norms and
$\|\mathfrak a_i(\mu_i s_i-\rho)\|_{E^*}\to0$. The definition of operator norm
then gives \eqref{eq:trace-sensor-conclusion}, with one tail uniform in $y$.
\end{proof}

The marked-vector assumptions are separate analytic inputs. In the selected
construction let $\mathsf q_j$ and $\mathsf w_j$ be the transported chosen and reference
responses, and set $m_j=b_j^{-1}\mathsf q_j-\mathsf w_j$. The scalar assembly in
Section~\ref{sec:ex-lattice-trace-weighted-estimates}, applied at the
postselected levels and under its stated response, price and scale inputs, gives
$m_{n_i}\to x$ and $\mathfrak a_i\rho(m_{n_i})\to0$. In particular, the
scale comparison above and $\widehat\beta_{n_i}\to\infty$ give
$b_{n_i}\ge1$ after discarding finitely many terms; the cited scalar
calculation supplies $\mathfrak a_i/b_{n_i}\to0$.
The distinct priced-mission argument in
Lemma~\ref{lem:ex-trace-small-means} gives the stronger field estimate
$\widehat\beta_{n_i}s_i(m_{n_i})\to0$, under its stated medium-core representation, oddness, response and scaling
identities, detector and local-price bounds, and scalar rate input. For large $t$, $\tfrac15(\log\log t)^2\le\log t$,
so $A_{1/5}(n_i)\le\widehat\beta_{n_i}$ eventually. Therefore
\[
 \|\mathfrak a_i s_i(m_{n_i})\|\le
 \|\widehat\beta_{n_i}s_i(m_{n_i})\|\longrightarrow0.
\]
Thus the same selected medium supplies all three marked-vector limits in
the proposition, subject to those two separate groups of inputs. The
finite-volume price, selected-chart and probe estimates remain draft
prerequisites. This proposition neither proves those producers nor asserts
the full existence theorem.
